\section{Conclusion}

This thesis investigated whether cryptocurrency price dynamics---Bitcoin in particular---contain statistically meaningful information for forecasting time-varying dependency structures with conventional asset markets, and whether such dependency forecasts can be operationalized as a practically relevant early-warning signal for investors. To address these questions, a reproducible two-layer analytical framework was developed and evaluated within a strict walk-forward experimental design. The first layer addressed the direct out-of-sample forecasting of rolling pairwise correlations between Bitcoin and a set of conventional assets. The second layer converted the output of the first stage into a binary stress-regime classification problem, generating an investor-facing warning signal for elevated downside risk.

The empirical evidence establishes that the intermarket dependency process is forecastable out of sample. Prediction accuracy is consistently positive across the evaluation set, with the strongest performance recorded for cross-cryptocurrency pairs and for medium-to-long rolling window specifications. The primary source of this predictability is, however, the strong serial persistence inherent in the rolling correlation target. Autoregressive and naive persistence baselines consequently dominate the average performance ranking across asset pairs. This finding does not undermine the thesis; rather, it precisely characterizes the structure of the forecasting problem and provides a methodologically honest account of what the data can support. It establishes that the dependency process exhibits memory, and that the dominant predictive signal is contained in the recent path of the correlation series itself.

Machine-learning models contribute meaningfully to the analysis, albeit in a comparative and structural role rather than as unconditional top performers. They frequently outperform the leakage-safe DCC-GARCH benchmark and provide a broader comparison class against which the limits of purely persistence-based prediction can be assessed. Their relative advantage is most pronounced in the second-layer signal problem, where the shift from regression to binary stress classification aligns more favorably with the strengths of ensemble methods. The investor signal layer demonstrates that crypto-derived features can be transformed into a statistically meaningful early-warning mechanism for conventional market stress regimes. The signal is moderate in magnitude rather than definitive, and its practical utility is best understood as a risk-management overlay---a supplementary input to exposure management and risk-limit setting---rather than as a standalone directional predictor.

The thesis thereby establishes a defensible academic contribution and a bounded but coherent practical contribution. Academically, it demonstrates that time-varying intermarket dependence between Bitcoin and conventional assets is forecastable in a rigorous out-of-sample framework, and it provides a systematic comparison of persistence baselines, machine-learning methods, and a classical econometric benchmark under consistent evaluation conditions. Practically, it shows that dependency forecasts informed by cryptocurrency dynamics can support an investor-facing warning mechanism, with the strongest and most robust signal performance observed for broad equity market indices.

Future extensions of this work should address the following priorities: testing the forecasting framework across multiple prediction horizons to assess whether machine-learning methods gain relative to persistence baselines as the forecast horizon lengthens; augmenting the feature set with macro-financial predictors such as implied volatility indices, credit spreads, term structure variables, and on-chain cryptocurrency metrics; replacing the rolling Pearson correlation target with richer dependence representations including copula-based measures and tail dependence coefficients; and integrating the signal layer into a portfolio optimization framework with explicit transaction cost and turnover constraints, enabling evaluation against utility-relevant performance criteria. Even within its current scope, this thesis provides a rigorous and substantively informative answer to the research problem and contributes to the evidence base for the role of cryptocurrency dynamics in forecasting and managing cross-market risk in conventional asset portfolios.
